%% 2i2d-familles-contrainte-deformation — allure qualitative de l'essai de traction
%% pour quatre familles de matériaux. Aucune valeur chiffrée : seules comptent la
%% pente (rigidité E) et la longueur de la partie plastique (ductilité).
\documentclass[tikz,border=4pt]{standalone}
\usepackage{liaisons}
\begin{document}
\begin{tikzpicture}[x=1cm,y=1cm,
  axe/.style={-{Stealth[length=5pt]}, line width=0.7pt, solideF},
  rappel/.style={line width=0.4pt, dash pattern=on 2.5pt off 2pt, black!55},
  cote/.style={{Stealth[length=4pt]}-{Stealth[length=4pt]}, line width=0.6pt, solideD},
  etiq/.style={font=\small, inner sep=2pt},
  petit/.style={font=\scriptsize, inner sep=1.5pt}]

%% --- les axes ---------------------------------------------------------------
\draw[axe] (0,0) -- (0,3.5);
\draw[axe] (0,0) -- (5.5,0);
\node[petit, anchor=north east] at (-0.02,-0.02) {$0$};
\node[petit, anchor=south east] at (5.32,0.07) {Déformation $\varepsilon = \Delta L / L$};
\node[petit, anchor=south west] at (-0.35,3.55) {Contrainte $\sigma = F/S$};

%% --- A : céramique / verre — fragile ------------------------------------------
\draw[line width=1.1pt, solideA] (0,0) -- (0.62,3.1);
\draw[line width=1.1pt, solideA] (0.5,2.98) -- (0.76,3.22) (0.5,3.22) -- (0.76,2.98);
\node[petit, anchor=west, text=solideA, align=left] at (0.82,3.2)
     {rupture \textbf{sans}\\déformation plastique};

%% --- B : acier doux — ductile ---------------------------------------------------
\draw[rappel] (0,1.9) -- (0.95,1.9);
\draw[rappel] (0,2.62) -- (3.15,2.62);
\node[petit, anchor=east] at (-0.03,1.9) {$R_e$};
\node[petit, anchor=east] at (-0.03,2.62) {$R_r$};
\draw[line width=1.6pt, solideB, line cap=round, line join=round] (0,0) -- (0.95,1.9)
  .. controls (1.02,2.06) and (1.12,1.72) .. (1.25,1.70)
  .. controls (1.45,1.68) and (1.55,1.78) .. (1.75,1.86)
  .. controls (2.35,2.25) and (2.70,2.54) .. (3.15,2.62)
  .. controls (3.5,2.67) and (3.72,2.5) .. (3.92,2.2);
\draw[line width=1pt, solideB] (3.80,2.32) -- (4.04,2.10) (3.80,2.10) -- (4.04,2.32);
\node[petit, anchor=west, text=solideB] at (4.05,2.22) {rupture};
\draw[line width=0.5pt, black!60] (0.34,0.68) -- (0.6,0.68) -- (0.6,1.2);
\node[petit, anchor=north west, inner sep=1pt] at (0.6,1.2) {$E$};

%% --- C : alliage d'aluminium ----------------------------------------------------
\draw[line width=1.1pt, solideE, line cap=round] (0,0) -- (0.85,1.05)
  .. controls (1.4,1.35) and (2.0,1.55) .. (2.6,1.62)
  .. controls (2.95,1.65) and (3.1,1.5) .. (3.2,1.3);
\draw[line width=0.9pt, solideE] (3.1,1.42) -- (3.3,1.22) (3.1,1.22) -- (3.3,1.42);

%% --- D : élastomère --------------------------------------------------------------
\draw[line width=1.1pt, solideD, line cap=round] (0,0)
  .. controls (1.6,0.28) and (3.4,0.42) .. (4.4,0.62)
  .. controls (4.95,0.78) and (5.2,1.25) .. (5.3,1.85);

%% --- légende ---------------------------------------------------------------------
\begin{scope}[shift={(4.88,3.42)}]
  \foreach \i/\c/\t in {0/solideA/{céramique (fragile)}, 1/solideB/{acier doux (ductile)},
                        2/solideE/{alliage d'aluminium}, 3/solideD/{élastomère}} {
    \draw[line width=1.2pt, \c] (0,-0.42*\i) -- (0.5,-0.42*\i);
    \node[petit, anchor=west, text=\c] at (0.58,-0.42*\i) {\t}; }
\end{scope}

%% --- zones élastique et plastique de l'acier ---------------------------------------
\draw[rappel] (0.95,1.9) -- (0.95,-0.7)  (3.92,2.2) -- (3.92,-0.7);
\draw[cote] (0.03,-0.55) -- (0.95,-0.55);
\draw[cote] (0.95,-0.55) -- (3.92,-0.55);
\node[petit, anchor=north, text=solideD] at (0.5,-0.6) {zone élastique};
\node[petit, anchor=north, text=solideD] at (2.4,-0.6) {zone plastique};

%% --- cartouche -------------------------------------------------------------------------
\node[draw=black!45, fill=black!3, rounded corners=3pt, inner sep=4pt, anchor=north,
      align=center, font=\scriptsize] at (3.7,-1.02)
     {la \textbf{pente} donne la rigidité $E$ ;\\
      la \textbf{longueur} de la partie plastique donne la ductilité};
\end{tikzpicture}
\end{document}
